\documentclass{zhvt-classic}

% 使用本地源流明体（GenRyuMin2 TW）按文件路径加载，不依赖系统字体名
% 版框（内框／界栏／书口侧栏／鱼尾）改用与版心印章同一印泥棕色：
%   印章图片 yshs.png 实色 = RGB(135,68,52) = HTML 874434
\usepackage{xcolor}
\usepackage{tikz}
% sealbrown 双注册：xcolor 供 TikZ/pgf，L3 供 cls 版框
\definecolor{sealbrown}{HTML}{874434}
\ExplSyntaxOn
\color_set:nnn { sealbrown } { HTML } { 874434 }
\ExplSyntaxOff

\zhvtset{
  grid color       = sealbrown,
  main font        = GenRyuMin2TW-SB,
  main font path   = C:/Users/Administrator/texmf/fonts/opentype/public/GenRyuMin2TW/,
jiazhu font      = GenRyuMin2TW-M,
  jiazhu font path = C:/Users/Administrator/texmf/fonts/opentype/public/GenRyuMin2TW/,
  % 页码框容纳 4 字：正常页「九十七」+前置空 1 格=4 行恰好；
  % 100+ 页「一百零一」共 5 行超出，末字「一/二/三」被框裁掉（2026-10-04 实测），故放宽到 5
  shukou page number chars = 5,
  % 书口书名框高从默认 6 字减到 4 字，把省出的 2 字高拨给 chapter 框，
  % 使 6 字书名+卷名放进 9 字高的 chapter 框后不再压线（2026-10-05）
  shukou title chars = 4,
}

\title{史记}[苏秦列传]
\author{汉·司马迁}
\maker{闲叟漫记·重製}


% ===== 索隐述赞样式（2026-11-08 四次定稿 · 愚大人亲手改定参数）=====
% ⚠ 竖排前提：自定义字体家族必须带 RawFeature={vertical:+vert}；
%   simhei.ttf 加 +vhal 会整体加载失败静默回退，只写 +vert；simfang 同理。
\newCJKfontfamily{\shuzanfont}[Path=C:/Windows/Fonts/,RawFeature={vertical:+vert}]{simhei.ttf}
\newCJKfontfamily{\shuzanbody}[Path=C:/Windows/Fonts/,RawFeature={vertical:+vert}]{simfang.ttf}
% 题头：152x30pt 圆角 6pt 棕底白字；baseline (0,2) 使整框落入界行格子不压版框线；
% 白字节点 x 不补偿、y +8.2pt（像素级墨迹实测 dx/dy≈0，L/R 留白约 20pt）。
\newcommand{\shuzanhead}[1]{%
  \begin{tikzpicture}[baseline={(0,2)}]
    \fill[sealbrown, rounded corners=6pt]
      (-76pt,-15pt) rectangle (76pt,15pt);
    \node[text=white, inner sep=0pt, anchor=center] at (0pt,8.2pt)
      {\shuzanfont\fontsize{23}{23}\selectfont #1};
  \end{tikzpicture}%
}
% ★ 修正书口卷名位置：chapter 框起始位置改回 cls 默认值，避免全局影响。
%   只在每个卷首页通过 \juan 命令里的空命令微调该页位置（安全，不会跨页污染）。
\ExplSyntaxOn
\AtBeginDocument{
  \dim_set:Nn \l__zhvt_shukou_chapter_title_start_dim
    { \l__zhvt_shukou_title_height_dim + \l__zhvt_font_size_dim * 2 }
  \dim_set:Nn \l__zhvt_shukou_page_number_start_dim
    { \l__zhvt_shukou_chapter_title_start_dim + \l__zhvt_shukou_chapter_title_height_dim }
}
\ExplSyntaxOff

\begin{document}

% ★★ 版心（书口）排布覆写：把书名放到鱼尾之下（2026-10-04）
%
%   cls 默认版心自上而下：书名（title 静态框，在鱼尾之上）→ 鱼尾 → 篇名（chapter 动态框）→ 页码
%   本模板改为：          鱼尾（位置不动）→ 书名 → 篇名 → 页码分割线 → 页码
%
%   手法：把 cls 放在 title 静态框里的书名，改放进 chapter 动态框（鱼尾之下）；
%         title 静态框留空。书名取 \l__zhvt_title_str，即 \title{...} 的第一参数，
%         所以换书时只改 \title，不必在此重复写书名。
%
%   ⚠ 三条必须遵守的写法（2026-10-04 实测）：
%     ① 汉字必须用 \hbox{...} 包住。flowfram 会把框内容装进一个宽 14.4mm（= 一个 line_sep
%        = 书口栏宽）的纵向 minipage 并按该宽度折行，而整页靠 PDF /Rotate 90 显示、
%        框内文字本身不变向 —— 不包 \hbox 时「围炉夜话」4 字会被逐字折成 4 行，
%        在版面上就成了横排 4 字并横向溢出版框。
%     ② 竖排下**一个 \hbox 就是一整列**。书名与篇名要同处一列，必须放进同一个 \hbox；
%        分成两个 \hbox 时，后一个会被当成新的一列、甩到隔壁栏去
%        （实测「序」跑到书口栏外侧 14.4mm 的正文列上）。中间的 \quad 在框内是横向空白，
%        显示出来正好是书名与篇名之间的纵向间隔。
%     ③ ★ 内容盒要写成 **与框等宽的定宽盒** `\hbox to \l__zhvt_shukou_chapter_title_height_dim{…\hfil}`。
%        实测：chapter 框是把内容盒沿纵向**居中**摆放的，内容宽度一变，书名就跟着漂：
%        自然宽度的 \hbox（书名 38.4mm）落在 x=157.2mm，加篇名后（57.6mm）落在 x=152.4mm，
%        目录页（无篇名）与正文页的书名就差半格。做成等宽盒后，无论有没有篇名，
%        内容固定从框顶起排，实测两页均落在 x=142.8mm，不再漂。
\ExplSyntaxOn
% 把书口 chapter 框（书名+篇名）起始位置下移 1 字高，避免 4 字卷名顶到鱼尾
% ⚠ cls 的 \__zhvt_calculate_dimensions: 在 \AtBeginDocument 里执行，会重算该 dim；
%   必须把覆写挂到 begin document 之后，否则会被二次覆盖（2026-10-05 实测）。
% (旧方案已移到导言区，直接修改 chapter 框起始位置)
% 卷标题命令 #1=书口篇名(鱼尾) #2=正文卷名（含卷序，卷名内空格用 ~）
%
% ★★ 目录页两项修订（2026-10-05，愚大人指出）★★
%   ① 页码起点统一对齐：
%      原状 —— 目录项写成 \par\hspace{4em}卷名\hspace{4em}页码，页码起点 = 卷名实宽 + 4em。
%      卷名实宽随字数变：「卷一~梁惠王上」6 字 +~、「卷七~离娄上」5 字 +~，
%      差 1 个字位 → 卷七~卷十 4 条的页码整体比其余 10 条高 27.21pt（一个字），看着上下不齐。
%      改法 —— 卷名补全角空格（宽 = 1 字 = 27.215pt）补足到 **7 字符位**再写进目录，
%      14 条页码起点遂全部落在同一条竖线上（起点对齐；末字随页码字数自然参差）。
%      ⚠ 只对目录里的副本 \l_tmpc_tl 补，正文卷名仍用 \l_tmpa_tl 原样 —— 正文像素不受影响。
%   ② 卷名整体上移半字位：前导 \hspace 由 4em 改 3.5em。
%      实测原状目录块在版框内偏下（天头留白 139.6pt / 地脚留白 110.7pt），
%      上移 0.5em = 13.6pt 后上下基本均分，视觉居中。
\NewDocumentCommand { \juan } { mm } {
	\clearpage
	\tl_set:Nn \l__zhvt_chapter_title_tl { #1 }
	\tl_set:Nn \l_tmpa_tl { #2 }
	\tl_set_eq:NN \l_tmpc_tl \l_tmpa_tl
	\int_step_inline:nn { 7 - \str_count:n { #2 } }
	  { \tl_put_right:Nn \l_tmpc_tl { 　 } }
	\tl_set:Nx \l_tmpb_tl { \zhnum { page } }
	\tl_put_left:Nn \l_tmpb_tl { \hspace { 4em } }
	\tl_put_left:Nn \l_tmpb_tl { \l_tmpc_tl }
	\tl_put_left:Nn \l_tmpb_tl { \par \hspace { 3.5em } }
	\addtocontents { toc } { \l_tmpb_tl }
	\skip_horizontal:n { \l__zhvt_font_size_dim * 4 }
	\l_tmpa_tl
	\par
	\thispagestyle { empty }
}
\cs_set_protected:Npn \__zhvt_set_shukou_content: {
	% ★ 鱼尾上方留白（title 静态框）嵌入闲叟钤印：
	%   整页 /Rotate 90，故此处把图用 \rotatebox{90} 转正为竖章；
	%   尺寸卡在版心条幅内（书内约 7mm 宽 × 33.8mm 高，留余量不压鱼尾）。
	\setstaticcontents * { title } {
		\vbox_to_ht:nn { \l__zhvt_baseline_skip_dim } {
			\vskip 14.2pt
			\hbox_to_wd:nn { \l__zhvt_shukou_title_height_dim } {
				\hfil
				\rotatebox [ origin = c ] { 90 } { \includegraphics [ width = 7mm , height = 33.8mm ] { yshs.png } }
				\hfil
			}
		}
	}
% ★ 版心（书口）统一改法（2026-11-16 愚大人定）：18pt＋框内居中＋书口文字同版框色
%   ① 原 27.2pt 时 7 字长篇名（如「廉颇蔺相如列传」）末字溢出 chapter 框、压住页码；
%      改 \fontsize{18pt}{18pt}，前后 \hfil 令字串在框内居中（不随篇名字数漂移）。
%   ② \vbox{\vskip 8.2pt} 定量下移，抵消缩字号后栏宽方向 8.2pt 上偏（像素实测 y 中心 543.3 = 原版）。
%      ⚠ \raisebox / \vspace* 在 flowfram 框内无效；换字号须同步换补偿：16pt→10.1、19pt→7.4。
%   ③ chapter 与 page~number 均置 \color{sealbrown}（HTML 874434），与版框／界行同色；正文仍黑。
%   ⚠ 页码字位、钤印、鱼尾、界行、正文一律未动；五部合订本同此口径。
	\setdynamiccontents * { chapter }
	{ \color{sealbrown} \vbox { \vskip 8.2pt \hbox to \l__zhvt_shukou_chapter_title_height_dim
		{ \hfil \fontsize{18pt}{18pt}\selectfont \l__zhvt_title_str \quad \l__zhvt_chapter_title_tl \hfil } } }
	\setdynamiccontents * { page~number } { \color{sealbrown} \quad \zhnum { page } }
}
\__zhvt_set_shukou_content:
\ExplSyntaxOff

\maketitle

% 单篇，无多卷目录，不渲染目录页

\juan{苏秦列传}{卷六十九~苏秦列传第九}
\hspace*{2em}苏秦者。东周雒阳人也。东事师於齐。而习之於鬼谷先生。\\
\hspace*{2em}出游数岁。大困而归。兄弟嫂妹妻妾窃皆笑之。曰。﹁周人之俗。治产业。力工商。逐什二以为务。今子释本而事口舌。困。不亦宜乎。﹂苏秦闻之而惭。自伤。乃闭室不出。出其书遍观之。曰。﹁夫士业已屈首受书。而不能以取尊荣。虽多亦奚以为。﹂於是得周书阴符。伏而读之。期年。以出揣摩。曰。﹁此可以说当世之君矣。﹂求说周显王。显王左右素习知苏秦。皆少之。弗信。\\
\hspace*{2em}乃西至秦。秦孝公卒。说惠王曰。﹁秦四塞之国。被山带渭。东有关河。西有汉中。南有巴蜀。北有代马。此天府也。以秦士民之众。兵法之教。可以吞天下。称帝而治。﹂秦王曰。﹁毛羽未成。不可以高蜚。文理未明。不可以并兼。﹂方诛商鞅。疾辩士。弗用。\\
\hspace*{2em}乃东之赵。赵肃侯令其弟成为相。号奉阳君。奉阳君弗说之。\\
\hspace*{2em}去游燕。岁馀而后得见。说燕文侯曰。燕东有朝鲜。辽东。北有林胡。楼烦。西有云中。九原。南有嘑沱。易水。地方二千馀里。带甲数十万。车六百乘。骑六千匹。粟支数年。南有碣石。雁门之饶。北有枣栗之利。民虽不佃作而足於枣栗矣。此所谓天府者也。\\
\hspace*{2em}﹁夫安乐无事。不见覆军杀将。无过燕者。大王知其所以然乎。夫燕之所以不犯寇被甲兵者。以赵之为蔽其南也。秦赵五战。秦再胜而赵三胜。秦赵相毙。而王以全燕制其后。此燕之所以不犯寇也。且夫秦之攻燕也。逾云中。九原。过代。上谷。弥地数千里。虽得燕城。秦计固不能守也。秦之不能害燕亦明矣。今赵之攻燕也。发号出令。不至十日而数十万之军军於东垣矣。渡呼沱。涉易水。不至四五日而距国都矣。故曰秦之攻燕也。战於千里之外。赵之攻燕也。战於百里之内。夫不忧百里之患而重千里之外。计无过於此者。是故原大王与赵从亲。天下为一。则燕国必无患矣。﹂\\
\hspace*{2em}文侯曰。﹁子言则可。然吾国小。西迫强赵。南近齐。齐。赵强国也。子必欲合从以安燕。寡人请以国从。﹂\\
\hspace*{2em}於是资苏秦车马金帛以至赵。而奉阳君已死。即因说赵肃侯曰。天下卿相人臣及布衣之士。皆高贤君之行义。皆原奉教陈忠於前之日久矣。虽然。奉阳君妒而君不任事。是以宾客游士莫敢自尽於前者。今奉阳君捐馆舍。君乃今复与士民相亲也。臣故敢进其愚虑。\\
\hspace*{2em}窃为君计者。莫若安民无事。且无庸有事於民也。安民之本。在於择交。择交而得则民安。择交而不得则民终身不安。请言外患。齐秦为两敌而民不得安。倚秦攻齐而民不得安。倚齐攻秦而民不得安。故夫谋人之主。伐人之国。常苦出辞断绝人之交也。原君慎勿出於口。请别白黑所以异。阴阳而已矣。君诚能听臣。燕必致旃裘狗马之地。齐必致鱼盐之海。楚必致橘柚之园。韩。魏。中山皆可使致汤沐之奉。而贵戚父兄皆可以受封侯。夫割地包利。五伯之所以覆军禽将而求也。封侯贵戚。汤武之所以放弑而争也。今君高拱而两有之。此臣之所以为君原也。\\
\hspace*{2em}今大王与秦。则秦必弱韩。魏。与齐。则齐必弱楚。魏。魏弱则割河外。韩弱则效宜阳。宜阳效则上郡绝。河外割则道不通。楚弱则无援。此三策者。不可不孰计也。\\
\hspace*{2em}夫秦下轵道。则南阳危。劫韩包周。则赵氏自操兵。据卫取卷。则齐必入朝秦。秦欲已得乎山东。则壁举兵而乡赵矣。秦甲渡河逾漳。据番吾。则兵必战於邯郸之下矣。此臣之所为君患也。\\
\hspace*{2em}当今之时。山东之建国莫彊於赵。赵地方二千馀里。带甲数十万。车千乘。骑万匹。粟支数年。西有常山。南有河漳。东有清河。北有燕国。燕固弱国。不足畏也。秦之所害於天下者莫如赵。然而秦不敢举兵伐赵者。何也。畏韩。魏之议其后也。然则韩。魏。赵之南蔽也。秦之攻韩。魏也。无有名山大川之限。稍蚕食之。傅国都而止。韩。魏不能支秦。必入臣於秦。秦无韩。魏之规。则祸必中於赵矣。此臣之所为君患也。\\
\hspace*{2em}臣闻尧无三夫之分。舜无咫尺之地。以有天下。禹无百人之聚。以王诸侯。汤武之士不过三千。车不过三百乘。卒不过三万。立为天子。诚得其道也。是故明主外料其敌之彊弱。内度其士卒贤不肖。不待两军相当而胜败存亡之机固已形於胸中矣。岂揜於众人之言而以冥冥决事哉。\\
\hspace*{2em}臣窃以天下之地图案之。诸侯之地五倍於秦。料度诸侯之卒十倍於秦。六国为一。并力西乡而攻秦。秦必破矣。今西面而事之。见臣於秦。夫破人之与破於人也。臣人之与臣於人也。岂可同日而论哉。\\
\hspace*{2em}夫衡人者。皆欲割诸侯之地以予秦。秦成。则高台榭。美宫室。听竽瑟之音。前有楼阙轩辕。后有长姣美人。国被秦患而不与其忧。是故夫衡人日夜务以秦权恐愒诸侯以求割地。故原大王孰计之也。\\
\hspace*{2em}﹁臣闻明主绝疑去谗。屏流言之迹。塞朋党之门。故尊主广地彊兵之计臣得陈忠於前矣。故窃为大王计。莫如一韩。魏。齐。楚。燕。赵以从亲。以畔秦。令天下之将相会於洹水之上。通质。刳白马而盟。要约曰。﹃秦攻楚。齐。魏各出锐师以佐之。韩绝其粮道。赵涉河漳。燕守常山之北。秦攻韩魏。则楚绝其后。齐出锐师而佐之。赵涉河漳。燕守云中。秦攻齐。则楚绝其后。韩守城皋。魏塞其道。赵涉河漳。博关。燕出锐师以佐之。秦攻燕。则赵守常山。楚军武关。齐涉勃海。韩。魏皆出锐师以佐之。秦攻赵。则韩军宜阳。楚军武关。魏军河外。齐涉清河。燕出锐师以佐之。诸侯有不如约者。以五国之兵共伐之。﹄六国从亲以宾秦。则秦甲必不敢出於函谷以害山东矣。如此。则霸王之业成矣。﹂\\
\hspace*{2em}赵王曰。﹁寡人年少。立国日浅。未尝得闻社稷之长计也。今上客有意存天下。安诸侯寡人敬以国从。﹂乃饰车百乘。黄金千溢。白璧百双。锦绣千纯。以约诸侯。\\
\hspace*{2em}是时周天子致文武之胙於秦惠王。惠王使犀首攻魏。禽将龙贾。取魏之雕阴。且欲东兵。苏秦恐秦兵之至赵也。乃激怒张仪。入之于秦。\\
\hspace*{2em}於是说韩宣王曰。韩北有巩。成皋之固。西有宜阳。商阪之塞。东有宛。穰。洧水。南有陉山。地方九百馀里。带甲数十万。天下之彊弓劲弩皆从韩出。谿子。少府时力。距来者。皆射六百步之外。韩卒超足而射。百发不暇止。远者括蔽洞胸。近者镝弇心。韩卒之剑戟皆出於冥山。棠谿。墨阳。合赙。邓师。宛冯。龙渊。太阿。皆陆断牛马。水截鹄雁。当敌则斩坚甲铁幕。革抉簠芮。无不毕具。以韩卒之勇。被坚甲。蹠劲弩。带利剑。一人当百。不足言也。夫以韩之劲与大王之贤。乃西面事秦。交臂而服。羞社稷而为天下笑。无大於此者矣。是故原大王孰计之。\\
\hspace*{2em}﹁大王事秦。秦必求宜阳。成皋。今兹效之。明年又复求割地。与则无地以给之。不与则弃前功而受后祸。且大王之地有尽而秦之求无已。以有尽之地而逆无已之求。此所谓市怨结祸者也。不战而地已削矣。臣闻鄙谚曰。﹃宁为鸡口。无为牛后。﹄今西面交臂而臣事秦。何异於牛后乎。夫以大王之贤。挟彊韩之兵。而有牛后之名。臣窃为大王羞之。﹂\\
\hspace*{2em}於是韩王勃然作色。攘臂瞋目。按剑仰天太息曰。﹁寡人虽不肖。必不能事秦。今主君诏以赵王之教。敬奉社稷以从。﹂\\
\hspace*{2em}又说魏襄王曰。大王之地。南有鸿沟。陈。汝南。许。郾。昆阳。召陵。舞阳。新都。新郪。东有淮。颍。枣。无胥。西有长城之界。北有河外。卷。衍。酸枣。地方千里。地名虽小。然而田舍庐庑之数。曾无所刍牧。人民之众。车马之多。日夜行不绝。輷輷殷殷。若有三军之众。臣窃量大王之国不下楚。然衡人怵王交彊虎狼之秦以侵天下。卒有秦患。不顾其祸。夫挟彊秦之势以内劫其主。罪无过此者。魏。天下之彊国也。王。天下之贤王也。今乃有意西面而事秦。称东籓。筑帝宫。受冠带。祠春秋。臣窃为大王耻之。\\
\hspace*{2em}臣闻越王勾践战敝卒三千人。禽夫差於干遂。武王卒三千人。革车三百乘。制纣於牧野。岂其士卒众哉。诚能奋其威也。今窃闻大王之卒。武士二十万。苍头二十万。奋击二十万。厮徒十万。车六百乘。骑五千匹。此其过越王勾践。武王远矣。今乃听於群臣之说而欲臣事秦。夫事秦必割地以效实。故兵未用而国已亏矣。凡群臣之言事秦者。皆奸人。非忠臣也。夫为人臣。割其主之地以求外交。偷取一时之功而不顾其后。破公家而成私门。外挟彊秦之势以内劫其主。以求割地。原大王孰察之。\\
\hspace*{2em}﹁周书曰。﹃緜緜不绝。蔓蔓柰何。豪氂不伐。将用斧柯。﹄前虑不定。后有大患。将柰之何。大王诚能听臣。六国从亲。专心并力壹意。则必无彊秦之患。故敝邑赵王使臣效愚计。奉明约。在大王之诏诏之。﹂\\
\hspace*{2em}魏王曰。﹁寡人不肖。未尝得闻明教。今主君以赵王之诏诏之。敬以国从。﹂\\
\hspace*{2em}因东说齐宣王曰。齐南有泰山。东有琅邪。西有清河。北有勃海。此所谓四塞之国也。齐地方二千馀里。带甲数十万。粟如丘山。三军之良。五家之兵。进如锋矢。战如雷霆。解如风雨。即有军役。未尝倍泰山。绝清河。涉勃海也。临菑之中七万户。臣窃度之。不下户三男子。三七二十一万。不待发於远县。而临菑之卒固已二十一万矣。临菑甚富而实。其民无不吹竽鼓瑟。弹琴击筑。斗鸡走狗。六博蹋鞠者。临菑之涂。车毂击。人肩摩。连衽成帷。举袂成幕。挥汗成雨。家殷人足。志高气扬。夫以大王之贤与齐之彊。天下莫能当。今乃西面而事秦。臣窃为大王羞之。\\
\hspace*{2em}且夫韩。魏之所以重畏秦者。为与秦接境壤界也。兵出而相当。不出十日而战胜存亡之机决矣。韩。魏战而胜秦。则兵半折。四境不守。战而不胜。则国已危亡随其后。是故韩。魏之所以重与秦战。而轻为之臣也。今秦之攻齐则不然。倍韩。魏之地。过卫阳晋之道。径乎亢父之险。车不得方轨。骑不得比行。百人守险。千人不敢过也。秦虽欲深入。则狼顾。恐韩。魏之议其后也。是故恫疑虚猲。骄矜而不敢进。则秦之不能害齐亦明矣。\\
\hspace*{2em}﹁夫不深料秦之无柰齐何。而欲西面而事之。是群臣之计过也。今无臣事秦之名而有彊国之实。臣是故原大王少留意计之。﹂\\
\hspace*{2em}齐王曰。﹁寡人不敏。僻远守海。穷道东境之国也。未尝得闻馀教。今足下以赵王诏诏之。敬以国从。﹂\\
\hspace*{2em}乃西南说楚威王曰。楚。天下之彊国也。王。天下之贤王也。西有黔中。巫郡。东有夏州。海阳。南有洞庭。苍梧。北有陉塞。郇阳。地方五千馀里。带甲百万。车千乘。骑万匹。粟支十年。此霸王之资也。夫以楚之彊与王之贤。天下莫能当也。今乃欲西面而事秦。则诸侯莫不西面而朝於章台之下矣。\\
\hspace*{2em}秦之所害莫如楚。楚彊则秦弱。秦彊则楚弱。其势不两立。故为大王计。莫如从亲以孤秦。大王不从。秦必起两军。一军出武关。一军下黔中。则鄢郢动矣。\\
\hspace*{2em}臣闻治之其未乱也。为之其未有也。患至而后忧之。则无及已。故原大王蚤孰计之。\\
\hspace*{2em}大王诚能听臣。臣请令山东之国奉四时之献。以承大王之明诏。委社稷。奉宗庙。练士厉兵。在大王之所用之。大王诚能用臣之愚计。则韩。魏。齐。燕。赵。卫之妙音美人必充后宫。燕。代橐驼良马必实外厩。故从合则楚王。衡成则秦帝。今释霸王之业。而有事人之名。臣窃为大王不取也。\\
\hspace*{2em}﹁夫秦。虎狼之国也。有吞天下之心。秦。天下之仇雠也。衡人皆欲割诸侯之地以事秦。此所谓养仇而奉雠者也。夫为人臣。割其主之地以外交彊虎狼之秦。以侵天下。卒有秦患。不顾其祸。夫外挟彊秦之威以内劫其主。以求割地。大逆不忠。无过此者。故从亲则诸侯割地以事楚。衡合则楚割地以事秦。此两策者相去远矣。二者大王何居焉。故敝邑赵王使臣效愚计。奉明约。在大王诏之。﹂\\
\hspace*{2em}楚王曰。﹁寡人之国西与秦接境。秦有举巴蜀并汉中之心。秦。虎狼之国。不可亲也。而韩。魏迫於秦患。不可与深谋。与深谋恐反人以入於秦。故谋未发而国已危矣。寡人自料以楚当秦。不见胜也。内与群臣谋。不足恃也。寡人卧不安席。食不甘味。心摇摇然如县旌而无所终薄。今主君欲一天下。收诸侯。存危国。寡人谨奉社稷以从。﹂\\
\hspace*{2em}於是六国从合而并力焉。苏秦为从约长。并相六国。\\
\hspace*{2em}北报赵王。乃行过雒阳。车骑辎重。诸侯各发使送之甚众。疑於王者。周显王闻之恐惧。除道。使人郊劳。苏秦之昆弟妻嫂侧目不敢仰视。俯伏侍取食。苏秦笑谓其嫂曰。﹁何前倨而后恭也。﹂嫂委地蒲服。以面掩地而谢曰。﹁见季子位高金多也。﹂苏秦喟然叹曰。﹁此一人之身。富贵则亲戚畏惧之。贫贱则轻易之。况众人乎。且使我有雒阳负郭田二顷。吾岂能佩六国相印乎。﹂於是散千金以赐宗族朋友。初。苏秦之燕。贷人百钱为资。乃得富贵。以百金偿之。遍报诸所尝见德者。其从者有一人独未得报。乃前自言。苏秦曰。我非忘子。子之与我至燕。再三欲去我易水之上。方是时。我困。故望子深。是以后子。子今亦得矣。苏秦既约六国从亲。归赵。赵肃侯封为武安君。乃投从约书於秦。秦兵不敢闚函谷关十五年。\\
\hspace*{2em}其后秦使犀首欺齐。魏。与共伐赵。欲败从约。齐。魏伐赵。赵王让苏秦。苏秦恐。请使燕。必报齐。苏秦去赵而从约皆解。\\
\hspace*{2em}秦惠王以其女为燕太子妇。是岁。文侯卒。太子立。是为燕易王。易王初立。齐宣王因燕丧伐燕。取十城。易王谓苏秦曰。﹁往日先生至燕。而先王资先生见赵。遂约六国从。今齐先伐赵。次至燕。以先生之故为天下笑。先生能为燕得侵地乎。﹂苏秦大惭。曰。﹁请为王取之。﹂\\
\hspace*{2em}苏秦见齐王。再拜。俯而庆。仰而吊。齐王曰。﹁是何庆吊相随之速也。﹂苏秦曰。﹁臣闻饥人所以饥而不食乌喙者。为其愈充腹而与饿死同患也。今燕虽弱小。即秦王之少婿也。大王利其十城而长与彊秦为仇。今使弱燕为雁行而彊秦敝其后。以招天下之精兵。是食乌喙之类也。﹂齐王愀然变色曰。﹁然则柰何。﹂苏秦曰。﹁臣闻古之善制事者。转祸为福。因败为功。大王诚能听臣计。即归燕之十城。燕无故而得十城。必喜。秦王知以己之故而归燕之十城。亦必喜。此所谓弃仇雠而得石交者也。夫燕。秦俱事齐。则大王号令天下。莫敢不听。是王以虚辞附秦。以十城取天下。此霸王之业也。﹂王曰。﹁善。﹂於是乃归燕之十城。\\
\hspace*{2em}人有毁苏秦者曰。﹁左右卖国反覆之臣也。将作乱。﹂苏秦恐得罪归。而燕王不复官也。苏秦见燕王曰。﹁臣。东周之鄙人也。无有分寸之功。而王亲拜之於庙而礼之於廷。今臣为王卻齐之兵而得十城。宜以益亲。今来而王不官臣者。人必有以不信伤臣於王者。臣之不信。王之福也。臣闻忠信者。所以自为也。进取者。所以为人也。且臣之说齐王。曾非欺之也。臣弃老母於东周。固去自为而行进取也。今有孝如曾参。廉如伯夷。信如尾生。得此三人者以事大王。何若。﹂王曰。﹁足矣。﹂苏秦曰。﹁孝如曾参。义不离其亲一宿於外。王又安能使之步行千里而事弱燕之危王哉。廉如伯夷。义不为孤竹君之嗣。不肯为武王臣。不受封侯而饿死首阳山下。有廉如此。王又安能使之步行千里而行进取於齐哉。信如尾生。与女子期於梁下。女子不来。水至不去。抱柱而死。有信如此。王又安能使之步行千里卻齐之彊兵哉。臣所谓以忠信得罪於上者也。﹂燕王曰。﹁若不忠信耳。岂有以忠信而得罪者乎。﹂苏秦曰。﹁不然。臣闻客有远为吏而其妻私於人者。其夫将来。其私者忧之。妻曰﹃勿忧。吾已作药酒待之矣﹄。居三日。其夫果至。妻使妾举药酒进之。妾欲言酒之有药。则恐其逐主母也。欲勿言乎。则恐其杀主父也。於是乎详僵而弃酒。主父大怒。笞之五十。故妾一僵而覆酒。上存主父。下存主母。然而不免於笞。恶在乎忠信之无罪也。夫臣之过。不幸而类是乎。﹂燕王曰。﹁先生复就故官。﹂益厚遇之。\\
\hspace*{2em}易王母。文侯夫人也。与苏秦私通。燕王知之。而事之加厚。苏秦恐诛。乃说燕王曰。﹁臣居燕不能使燕重。而在齐则燕必重。﹂燕王曰。﹁唯先生之所为。﹂於是苏秦详为得罪於燕而亡走齐。齐宣王以为客卿。\\
\hspace*{2em}齐宣王卒。湣王即位。说湣王厚葬以明孝。高宫室大苑囿以明得意。欲破敝齐而为燕。燕易王卒。燕哙立为王。其后齐大夫多与苏秦争宠者。而使人刺苏秦。不死。殊而走。齐王使人求贼。不得。苏秦且死。乃谓齐王曰。﹁臣即死。车裂臣以徇於市。曰﹃苏秦为燕作乱於齐﹄。如此则臣之贼必得矣。﹂於是如其言。而杀苏秦者果自出。齐王因而诛之。燕闻之曰。﹁甚矣。齐之为苏生报仇也。﹂\\
\hspace*{2em}苏秦既死。其事大泄。齐后闻之。乃恨怒燕。燕甚恐。苏秦之弟曰代。代弟苏厉。见兄遂。亦皆学。及苏秦死。代乃求见燕王。欲袭故事。曰。﹁臣。东周之鄙人也。窃闻大王义甚高。鄙人不敏。释鉏耨而干大王。至於邯郸。所见者绌於所闻於东周。臣窃负其志。及至燕廷。观王之群臣下吏。王。天下之明王也。﹂燕王曰。﹁子所谓明王者何如也。﹂对曰。﹁臣闻明王务闻其过。不欲闻其善。臣请谒王之过。夫齐。赵者。燕之仇雠也。楚。魏者。燕之援国也。今王奉仇雠以伐援国。非所以利燕也。王自虑之。此则计过。无以闻者。非忠臣也。﹂王曰。﹁夫齐者固寡人之雠。所欲伐也。直患国敝力不足也。子能以燕伐齐。则寡人举国委子。﹂对曰。﹁凡天下战国七。燕处弱焉。独战则不能。有所附则无不重。南附楚。楚重。西附秦。秦重。中附韩。魏。韩。魏重。且苟所附之国重。此必使王重矣。今夫齐。长主而自用也。南攻楚五年。畜聚竭。西困秦三年。士卒罢敝。北与燕人战。覆三军。得二将。然而以其馀兵南面举五千乘之大宋。而包十二诸侯。此其君欲得。其民力竭。恶足取乎。且臣闻之。数战则民劳。久师则兵敝矣。﹂燕王曰。﹁吾闻齐有清济。浊河可以为固。长城。钜防足以为塞。诚有之乎。﹂对曰。﹁天时不与。虽有清济。浊河。恶足以为固。民力罢敝。虽有长城。钜防。恶足以为塞。且异日济西不师。所以备赵也。河北不师。所以备燕也。今济西河北尽已役矣。封内敝矣。夫骄君必好利。而亡国之臣必贪於财。王诚能无羞从子母弟以为质。宝珠玉帛以事左右。彼将有德燕而轻亡宋。则齐可亡已。﹂燕王曰。﹁吾终以子受命於天矣。﹂燕乃使一子质於齐。而苏厉因燕质子而求见齐王。齐王怨苏秦。欲囚苏厉。燕质子为谢。已遂委质为齐臣。\\
\hspace*{2em}燕相子之与苏代婚。而欲得燕权。乃使苏代侍质子於齐。齐使代报燕。燕王哙问曰。﹁齐王其霸乎。﹂曰。﹁不能。﹂曰。﹁何也。﹂曰。﹁不信其臣。﹂於是燕王专任子之。已而让位。燕大乱。齐伐燕。杀王哙。子之。燕立昭王。而苏代。苏厉遂不敢入燕。皆终归齐。齐善待之。\\
\hspace*{2em}苏代过魏。魏为燕执代。齐使人谓魏王曰。﹁齐请以宋地封泾阳君。秦必不受。秦非不利有齐而得宋地也。不信齐王与苏子也。今齐魏不和如此其甚。则齐不欺秦。秦信齐。齐秦合。泾阳君有宋地。非魏之利也。故王不如东苏子。秦必疑齐而不信苏子矣。齐秦不合。天下无变。伐齐之形成矣。﹂於是出苏代。代之宋。宋善待之。\\
\hspace*{2em}齐伐宋。宋急。苏代乃遗燕昭王书曰。\\
\hspace*{2em}夫列在万乘而寄质於齐。名卑而权轻。奉万乘助齐伐宋。民劳而实费。夫破宋。残楚淮北。肥大齐。雠彊而国害。此三者皆国之大败也。然且王行之者。将以取信於齐也。齐加不信於王。而忌燕愈甚。是王之计过矣。夫以宋加之淮北。强万乘之国也。而齐并之。是益一齐也。北夷方七百里。加之以鲁。卫。彊万乘之国也。而齐并之。是益二齐也。夫一齐之彊。燕犹狼顾而不能支。今以三齐临燕。其祸必大矣。\\
\hspace*{2em}虽然。智者举事。因祸为福。转败为功。齐紫。败素也。而贾十倍。越王勾践栖於会稽。复残彊吴而霸天下。此皆因祸为福。转败为功者也。\\
\hspace*{2em}今王若欲因祸为福。转败为功。则莫若挑霸齐而尊之。使使盟於周室。焚秦符。曰﹁其大上计。破秦。其次。必长宾之﹂。秦挟宾以待破。秦王必患之。秦五世伐诸侯。今为齐下。秦王之志苟得穷齐。不惮以国为功。然则王何不使辩士以此言说秦王曰。﹁燕。赵破宋肥齐。尊之为之下者。燕。赵非利之也。燕。赵不利而势为之者。以不信秦王也。然则王何不使可信者接收燕。赵。令泾阳君。高陵君先於燕。赵。秦有变。因以为质。则燕。赵信秦。秦为西帝。燕为北帝。赵为中帝。立三帝以令於天下。韩。魏不听则秦伐之。齐不听则燕。赵伐之。天下孰敢不听。天下服听。因驱韩。魏以伐齐。曰﹃必反宋地。归楚淮北﹄。反宋地。归楚淮北。燕。赵之所利也。并立三帝。燕。赵之所原也。夫实得所利。尊得所原。燕。赵弃齐如脱鵕矣。今不收燕。赵。齐霸必成。诸侯赞齐而王不从。是国伐也。诸侯赞齐而王从之。是名卑也。今收燕。赵。国安而名尊。不收燕。赵。国危而名卑。夫去尊安而取危卑。智者不为也。﹂秦王闻若说。必若刺心然。则王何不使辩士以此若言说秦。秦必取。齐必伐矣。\\
\hspace*{2em}夫取秦。厚交也。伐齐。正利也。尊厚交。务正利。圣王之事也。\\
\hspace*{2em}燕昭王善其书。曰。﹁先人尝有德苏氏。子之之乱而苏氏去燕。燕欲报仇於齐。非苏氏莫可。﹂乃召苏代。复善待之。与谋伐齐。竟破齐。湣王出走。\\
\hspace*{2em}久之。秦召燕王。燕王欲往。苏代约燕王曰。楚得枳而国亡。齐得宋而国亡。齐。楚不得以有枳。宋而事秦者。何也。则有功者。秦之深雠也。秦取天下。非行义也。暴也。秦之行暴。正告天下。\\
\hspace*{2em}告楚曰。﹃蜀地之甲。乘船浮於汶。乘夏水而下江。五日而至郢。汉中之甲。乘船出於巴。乘夏水而下汉。四日而至五渚。寡人积甲宛东下随。智者不及谋。勇土不及怒。寡人如射隼矣。王乃欲待天下之攻函谷。不亦远乎。﹄楚王为是故。十七年事秦。\\
\hspace*{2em}秦正告韩曰。﹃我起乎少曲。一日而断大行。我起乎宜阳而触平阳。二日而莫不尽繇。我离两周而触郑。五日而国举。﹄韩氏以为然。故事秦。\\
\hspace*{2em}秦正告魏曰。﹃我举安邑。塞女戟。韩氏太原卷。我下轵。道南阳。封冀。包两周。乘夏水。浮轻舟。彊弩在前。錟戈在后。决荥口。魏无大梁。决白马之口。魏无外黄。济阳。决宿胥之口。魏无虚。顿丘。陆攻则击河内。水攻则灭大梁。﹄魏氏以为然。故事秦。\\
\hspace*{2em}秦欲攻安邑。恐齐救之。则以宋委於齐。曰。﹃宋王无道。为木人以寡人。射其面。寡人地绝兵远。不能攻也。王苟能破宋有之。寡人如自得之。﹄已得安邑。塞女戟。因以破宋为齐罪。\\
\hspace*{2em}秦欲攻韩。恐天下救之。则以齐委於天下。曰。﹃齐王四与寡人约。四欺寡人。必率天下以攻寡人者三。有齐无秦。有秦无齐。必伐之。必亡之。﹄已得宜阳。少曲。致蔺。石。因以破齐为天下罪。\\
\hspace*{2em}秦欲攻魏重楚。则以南阳委於楚。曰。﹃寡人固与韩且绝矣。残均陵。塞鄳阸戹。苟利於楚。寡人如自有之。﹄魏弃与国而合於秦。因以塞鄳阸戹为楚罪。\\
\hspace*{2em}兵困於林中。重燕。赵。以胶东委於燕。以济西委於赵。已得讲於魏。至公子延。因犀首属行而攻赵。\\
\hspace*{2em}兵伤於谯石。而遇败於阳马。而重魏。则以叶。蔡委於魏。已得讲於赵。则劫魏。不为割。困则使太后弟穰侯为和。嬴则兼欺舅与母。\\
\hspace*{2em}適燕者曰﹃以胶东﹄。適赵者曰﹃以济西﹄。適魏者曰﹃以叶。蔡﹄。適楚者曰﹃以塞鄳阸戹﹄。適齐者曰﹃以宋﹄。此必令言如循环。用兵如刺蜚。母不能制。舅不能约。﹁龙贾之战。岸门之战。封陵之战。高商之战。赵庄之战。秦之所杀三晋之民数百万。今其生者皆死秦之孤也。西河之外。上雒之地。三川晋国之祸。三晋之半。秦祸如此其大也。而燕。赵之秦者。皆以争事秦说其主。此臣之所大患也。﹂\\
\hspace*{2em}燕昭王不行。苏代复重於燕。\\
\hspace*{2em}燕使约诸侯从亲如苏秦时。或从或不。而天下由此宗苏氏之从约。代。厉皆以寿死。名显诸侯。\\
\hspace*{2em}太史公曰。苏秦兄弟三人。皆游说诸侯以显名。其术长於权变。而苏秦被反间以死。天下共笑之。讳学其术。然世言苏秦多异。异时事有类之者皆附之苏秦。夫苏秦起闾阎。连六国从亲。此其智有过人者。吾故列其行事。次其时序。毋令独蒙恶声焉。\\

% ===== 索隐述赞（司马贞《史记索隐》）· 统一样式 =====
% 题头「索隐述赞曰」五字并入反白框；正文每四句一列，定宽盒 \hfil 居列中；
% \fontsize{24}{41} 列距对齐全书界行（121+41k）；\clearpage 使述赞独占一页。
\clearpage
\shuzanhead{索隐述赞曰}\linebreak
{\shuzanbody\fontsize{24}{41}\selectfont
\hbox to \hsize{\hfil 季子周人。师事鬼谷。揣摩既就。阴符伏读。\hfil}\linebreak
\hbox to \hsize{\hfil 合从离衡。佩印者六。天王除道。家人扶服。\hfil}\linebreak
\hbox to \hsize{\hfil 贤哉代。厉。继荣党族。\hfil}%
\par}
\end{document}
